\documentclass[conference]{IEEEtran}
\usepackage{cite}
\usepackage{amsmath,amssymb,amsfonts}
\usepackage{graphicx}

\begin{document}

\title{COMP0245 Assignment}

\author{\IEEEauthorblockN{[Your Name], Ricktho Sarkar, Hafizh Ilham}
\IEEEauthorblockA{University College London\\
London, United Kingdom}
}

\maketitle

\begin{abstract}
\textit{[Summarise the problem, methodology, main experimental results, and conclusions once the report is complete.]}
\end{abstract}

\begin{IEEEkeywords}
mobile manipulation, robot control, linear regression
\end{IEEEkeywords}

\section{Introduction}
\subsection{Problem Formulation}
This project investigates control and learning for a simulated Robotnik mobile manipulator reaching a target position.
\textit{[Define the task, its inputs and outputs, and the criteria used to judge success.]}

\subsection{Assumptions and Constraints}
\textit{[State assumptions about the simulation, available state information, and target positions. Describe relevant constraints such as joint limits, actuator limits, workspace reachability, and computation or data budgets.]}

\section{Methodology}
\subsection{Implementation}
\textit{[Describe the simulation setup, mobile base and arm controllers, data collection, and linear regression training. Explain the data flow, model inputs and outputs, and training and evaluation procedure. Include equations or a system diagram where useful.]}

\subsection{Design Choices and Justification}
\textit{[Justify the controller design, model and feature choices, data collection strategy, and parameter settings. Explain the alternatives considered and relevant trade-offs.]}

\section{Results}
\subsection{Experimental Design and Comparisons}
\textit{[Define the experiments, target configurations, baselines, and evaluation metrics. Describe the training and test split, repeated trials, and conditions kept consistent across comparisons.]}

\subsection{Experimental Results}
\textit{[Report measured results in tables and plots. Suitable examples include end-effector error over time, reaching success rate, and prediction errors for arm torques and wheel velocities. Label axes and units, distinguish training from test results, and report variability where available.]}

\subsection{Discussion and Interpretation}
\textit{[Interpret the results and compare methods against the stated success criteria. Explain observed trends, failure cases, and differences between prediction accuracy and robot task performance.]}

\section{Conclusion}
\textit{[Summarise the main findings and explain how far the work meets the project objectives.]}

\subsection{Limitations}
\textit{[Discuss limitations of the method and experiments, including modelling assumptions, data coverage, generalisation, and the scope of the evaluation.]}

\subsection{Improvements and Future Work}
\textit{[Explain how the method and experiments could be improved. Relate each proposed improvement to a limitation or result, and describe how its benefit could be tested.]}

% Add a bibliography here when sources have been cited in the report.

\end{document}
